
\chapter{2016年浙江卷（理）}

\section{单选题}

\begin{problem}
已知集合 \( P=\{x\in\R\mid 1\le x\le 3\} \)，\( Q=\{x\in\R\mid x^2\ge 4\} \)，则 \( P\cup \left( \complement_{\R} Q \right)= \)（\quad）

\choices
    {\([2,3]\)}
    {\(( -2,3]\)}
    {\([1,2)\)}
    {\(( -\infty,-2]\cup [1,+\infty)\)}
\end{problem}

\begin{answer}
B
\end{answer}

\begin{solution}
由 \(x^2\geqslant4\) 得 \(Q=(-\infty,-2]\cup[2,+\infty)\)，所以
\(\complement_{\R}Q=(-2,2)\). 因而
\[
P\cup\left(\complement_{\R}Q\right)=[1,3]\cup(-2,2)=(-2,3]
\]
故选 B.
\end{solution}

\begin{problem}
已知互相垂直的平面 \(\alpha\)，\(\beta\) 交于直线 \( l \)，若直线 \(m\)，\(n\) 满足 \( m\parallel \alpha \)，\( n\perp \beta \)，则（\quad）

\choices
    {\(m\parallel l\)}
    {\(m\parallel n\)}
    {\(n\perp l\)}
    {\(m\perp n\)}
\end{problem}

\begin{answer}
C
\end{answer}

\begin{solution}
因为 \(n\perp\beta\)，所以直线 \(n\) 的方向垂直于平面 \(\beta\) 内任意直线的方向. 又 \(l\subset\beta\)，故 \(n\perp l\). 条件 \(m\parallel\alpha\) 不能确定 \(m\) 与 \(l\)、\(n\) 的关系，因此只有 C 正确.
\end{solution}

\begin{problem}
在平面上，过点 \(P\) 作直线 \( l \) 的垂线所得的垂足称为点 \(P\) 在直线 \( l \) 上的投影，由区域 \( \begin{cases}
x-2\le 0,\\
x+y\ge 0,\\
x-3y+4\ge 0
\end{cases} \) 中的点在直线 \( x+y-2=0 \) 上的投影构成的线段记为 \(AB\)，则 \( |AB|= \)（\quad）

\choices
    {\(2\sqrt2\)}
    {\(4\)}
    {\(3\sqrt2\)}
    {\(6\)}
\end{problem}

\begin{answer}
C
\end{answer}

\begin{solution}
区域的边界为 \(x=2\)、\(y=-x\) 和 \(y=\dfrac{x+4}{3}\)，其三个顶点为
\((-1,1)\)、\((2,-2)\)、\((2,2)\). 设原区域内点 \((x,y)\) 在直线 \(x+y=2\) 上的投影为
\(H=(1+t,1-t)\)，则投影方向为 \((1,1)\)，从而
\[
t=\frac{x-y}{2}
\]
在线性区域上，\(x-y\) 的最大值和最小值分别在顶点处取得：
\(\max(x-y)=4\)，\(\min(x-y)=-2\).
所以 \(t\) 的范围为 \([-1,2]\)，投影线段的长度为
\[
(2-(-1))\sqrt{1^2+(-1)^2}=3\sqrt2
\]
故选 C.
\end{solution}

\begin{problem}
命题“\( \forall x\in\R \)，\( \exists n\in\N^* \)，使得 \( n\ge x^2 \)”的否定形式是（\quad）

\choices
    {\( \forall x\in\R \)，\( \exists n\in\N^* \)，使得 \( n<x^2 \)}
    {\( \forall x\in\R \)，\( \forall n\in\N^* \)，使得 \( n<x^2 \)}
    {\( \exists x\in\R \)，\( \exists n\in\N^* \)，使得 \( n<x^2 \)}
    {\( \exists x\in\R \)，\( \forall n\in\N^* \)，使得 \( n<x^2 \)}
\end{problem}

\begin{answer}
D
\end{answer}

\begin{solution}
利用量词否定公式
\(\neg(\forall x\,\exists n\,P(x,n))\iff\exists x\,\forall n\,\neg P(x,n)\)，并且
\(\neg(n\geqslant x^2)\) 等价于 \(n<x^2\). 因此原命题的否定为
\(\exists x\in\R\)，对任意 \(n\in\N^*\) 均有 \(n<x^2\).
故选 D.
\end{solution}

\begin{problem}
设函数 \( f(x)=\sin^2x+b\sin x+c \)，则 \( f(x) \) 的最小正周期（\quad）

\choices
    {\text{与 \( b \) 有关，且与 \( c \) 有关}}
    {\text{与 \( b \) 有关，但与 \( c \) 无关}}
    {\text{与 \( b \) 无关，且与 \( c \) 无关}}
    {\text{与 \( b \) 无关，但与 \( c \) 有关}}
\end{problem}

\begin{answer}
B
\end{answer}

\begin{solution}
利用降幂公式，
\[
f(x)=\frac{1-\cos2x}{2}+b\sin x+c
\]
当 \(b=0\) 时，\(f(x)\) 的最小正周期为 \(\pi\)；当 \(b\ne0\) 时，含有 \(b\sin x\) 项，周期不能为 \(\pi\)，而 \(2\pi\) 是周期，所以最小正周期为 \(2\pi\). 常数 \(c\) 不影响周期，故周期与 \(b\) 有关但与 \(c\) 无关.
故选 B.
\end{solution}

\begin{problem}
如图，点列 \( \{A_n\}\)，\( \{B_n\} \) 分别在某锐角的两边上，且 \( |A_nA_{n+1}|=|A_{n+1}A_{n+2}| \)，\( A_n\ne A_{n+1} \)，\( n\in\N^* \)，\( |B_nB_{n+1}|=|B_{n+1}B_{n+2}| \)，\( B_n\ne B_{n+1} \)，\( n\in\N^* \)（\( P\ne Q \) 表示点 \(P\) 与点 \(Q\) 不重合）. 若 \( d_n=|A_nB_n| \)，\( S_n \) 为 \( \triangle A_nB_nB_{n+1} \) 的面积，则（\quad）

\bitmapfigure[max width=0.62\textwidth]{img/2016/zhejiang_science/q6_fig1.png}

\choices
    {\(\{S_n\}\) 是等差数列}
    {\(\{S_n^2\}\) 是等差数列}
    {\(\{d_n\}\) 是等差数列}
    {\(\{d_n^2\}\) 是等差数列}
\end{problem}

\begin{answer}
A
\end{answer}

\begin{solution}
设两边的夹角为 \(\theta\)，并令 \(OA_n=x_n\)，\(OB_n=y_n\). 由题图中点列沿两边依次排列，且相邻线段长度分别保持不变，设
\(A_nA_{n+1}=p\)，\(B_nB_{n+1}=q\)，则
\(x_n=x_1+(n-1)p\)，\(y_n=y_1+(n-1)q\).
三角形 \(A_nB_nB_{n+1}\) 以 \(B_nB_{n+1}\) 为底，其高为 \(x_n\sin\theta\)，因此
\[
S_n=\frac12q x_n\sin\theta
 =\frac12q\sin\theta\bigl(x_1+(n-1)p\bigr)
\]
这是关于 \(n\) 的一次式，所以 \(\{S_n\}\) 是等差数列. 故选 A.
\end{solution}

\begin{problem}
已知椭圆 \( C_1:\dfrac{x^2}{m^2}+y^2=1 \)（\( m>1 \)）与双曲线 \( C_2:\dfrac{x^2}{n^2}-y^2=1 \)（\( n>0 \)）的焦点重合，\(e_1\)，\(e_2\) 分别为 \(C_1\)，\(C_2\) 的离心率，则（\quad）

\choices
    {\(m>n\) 且 \( e_1e_2>1 \)}
    {\(m>n\) 且 \( e_1e_2<1 \)}
    {\(m<n\) 且 \( e_1e_2>1 \)}
    {\(m<n\) 且 \( e_1e_2<1 \)}
\end{problem}

\begin{answer}
A
\end{answer}

\begin{solution}
设两曲线的公共焦半距为 \(c\). 由焦点重合，
\(c^2=m^2-1=n^2+1\)，\(m^2-n^2=2\)
所以 \(m>n\). 又
\[
e_1e_2=\frac{c}{m}\cdot\frac{c}{n}=\frac{c^2}{mn}
\]
并且
\[
(c^2)^2-m^2n^2=(m^2-1)^2-m^2(m^2-2)=1>0
\]
由于 \(c^2\)，\(mn>0\)，得 \(c^2>mn\)，故 \(e_1e_2>1\). 因此选 A.
\end{solution}

\begin{problem}
已知实数 \(a\)，\(b\)，\(c\)，则（\quad）

\choices
    {\(\bigl|a^2+b+c\bigr|+\bigl|a+b^2+c\bigr|\le 1\) 时，\( a^2+b^2+c^2<100 \)}
    {\(\bigl|a^2+b+c\bigr|+\bigl|a^2+b-c\bigr|\le 1\) 时，\( a^2+b^2+c^2<100 \)}
    {\(\bigl|a+b+c^2\bigr|+\bigl|a+b-c^2\bigr|\le 1\) 时，\( a^2+b^2+c^2<100 \)}
    {\(\bigl|a^2+b+c\bigr|+\bigl|a+b^2-c\bigr|\le 1\) 时，\( a^2+b^2+c^2<100 \)}
\end{problem}

\begin{answer}
D
\end{answer}

\begin{solution}
先说明 A、B、C 均不成立.

选项 A 取 \(a=b=10\)，\(c=-110\)，则
\[
|a^2+b+c|+|a+b^2+c|=0
\]
但 \(a^2+b^2+c^2=12300>100\). 选项 B 取 \(a=10\)，\(b=-100\)，\(c=0\)，则两个绝对值也都为 \(0\)，而 \(a^2+b^2+c^2=10100>100\). 选项 C 取 \(a=10\)，\(b=-10\)，\(c=0\)，同样两个绝对值都为 \(0\)，而 \(a^2+b^2+c^2=200>100\).

对于选项 D，令
\(u=a^2+b+c\)，\(v=a+b^2-c\).
由条件 \(|u|+|v|\leqslant1\) 得
\[
|a^2+a+b^2+b|=|u+v|\leqslant1
\]
于是
\[
\left(a+\frac12\right)^2+\left(b+\frac12\right)^2\leqslant\frac32
\]
从而 \(|a|<2\)，\(|b|<2\). 另外，
\[
|a^2-a+b-b^2+2c|=|u-v|\leqslant1
\]
由 \(|a^2-a|<6\)、\(|b-b^2|<6\)，得
\(2|c|<1+6+6=13\)，\(|c|<\frac{13}{2}\).
因此
\[
a^2+b^2+c^2<4+4+\left(\frac{13}{2}\right)^2=\frac{201}{4}<100
\]
所以 D 正确，故选 D.
\end{solution}

\section{填空题}

\begin{problem}
若抛物线 \( y^2=4x \) 上的点 \(M\) 到焦点的距离为 \( 10 \)，则 \(M\) 到 \( y \) 轴的距离是\fillinblank{}.
\end{problem}

\begin{answer}
9
\end{answer}

\begin{solution}
抛物线 \(y^2=4x\) 的焦点为 \(F(1,0)\)，准线为 \(x=-1\). 由抛物线的定义，点 \(M\) 到焦点的距离等于点 \(M\) 到准线的距离，故
\(x_M+1=10\)，\(x_M=9\).
因为抛物线在 \(x\geqslant0\) 上，所以 \(M\) 到 \(y\) 轴的距离为 \(x_M=9\).
\end{solution}

\begin{problem}
已知 \( 2\cos^2x+\sin 2x=A\sin(\omega x+\varphi)+b \)（\( A>0 \)），则 \( A=\fillinblank{} \)，\( b=\fillinblank{} \).
\end{problem}

\begin{answer}
\(\sqrt2\)，\(1\)
\end{answer}

\begin{solution}
由降幂公式，
\[
2\cos^2x+\sin2x=1+\cos2x+\sin2x
 =1+\sqrt2\sin\left(2x+\frac\pi4\right)
\]
与 \(A\sin(\omega x+\varphi)+b\) 比较，且 \(A>0\)，可得 \(A=\sqrt2\)，\(b=1\).
\end{solution}

\begin{problem}
某几何体的三视图如图所示（单位：\(\mathrm{cm}\)），则该几何体的表面积是\fillinblank{} \(\mathrm{cm}^2\)，体积是\fillinblank{} \(\mathrm{cm}^3\).

\bitmapfigure[max width=0.42\textwidth]{img/2016/zhejiang_science/q11_fig1.png}
\end{problem}

\begin{answer}
\(72\)，\(32\)
\end{answer}

\begin{solution}
由三视图可还原出该几何体由两个全等的长方体组成，每个长方体的长、宽、高分别为 \(4\mathrm{~cm}\)、\(2\mathrm{~cm}\)、\(2\mathrm{~cm}\)，两个长方体共用一个面积为 \(2\times2=4\mathrm{~cm}^2\) 的面. 因此体积为
\[
V=2\times(4\times2\times2)=32\mathrm{~cm}^3
\]
每个长方体的表面积为
\[
S_0=2(4\times2+4\times2+2\times2)=40\mathrm{~cm}^2
\]
共用的面在两个长方体的表面积中各计算了一次，但它不是外表面，所以该几何体的表面积为
\[
S=2S_0-2\times4=72\mathrm{~cm}^2
\]
\end{solution}

\begin{problem}
已知 \( a>b>1 \)，若 \( \log_a b+\log_b a=\dfrac52 \)，\( a^b=b^a \)，则 \( a=\fillinblank{} \)，\( b=\fillinblank{} \).
\end{problem}

\begin{answer}
\(4\)，\(2\)
\end{answer}

\begin{solution}
令 \(t=\log_a b\). 因为 \(a>b>1\)，所以 \(0<t<1\)，且 \(\log_ba=\dfrac1t\). 由已知，
\(t+\frac1t=\frac52\)，\(\Longrightarrow\)，\(2t^2-5t+2=0\).
解得 \(t=2\) 或 \(t=\dfrac12\)，结合 \(0<t<1\) 得 \(t=\dfrac12\)，故 \(b=a^{1/2}\)，即 \(a=b^2\).

对 \(a^b=b^a\) 两边取自然对数，得 \(b\ln a=a\ln b\). 代入 \(a=b^2\)，并利用 \(b>1\)，得
\(2b\ln b=b^2\ln b\)，\(\Longrightarrow\)，\(b=2\).
所以 \(a=4\)，\(b=2\).
\end{solution}

\begin{problem}
设数列 \( \{a_n\} \) 的前 \( n \) 项和为 \( S_n \)，若 \( S_2=4 \)，\( a_{n+1}=2S_n+1 \)，\( n\in\N^* \)，则 \( a_1=\fillinblank{} \)，\( S_5=\fillinblank{} \).
\end{problem}

\begin{answer}
\(1\)，\(121\)
\end{answer}

\begin{solution}
由 \(S_2=a_1+a_2=4\)，以及 \(n=1\) 时的递推式 \(a_2=2S_1+1=2a_1+1\)，得
\[
a_1+(2a_1+1)=4
\]
所以 \(a_1=1\)，\(a_2=3\). 对 \(n\geqslant2\)，将相邻两式相减，得
\[
a_{n+1}-a_n=(2S_n+1)-(2S_{n-1}+1)=2a_n
\]
即 \(a_{n+1}=3a_n\). 结合 \(a_1=1\)，得 \(a_n=3^{n-1}\). 因此
\[
S_5=1+3+9+27+81=121
\]
\end{solution}

\begin{problem}
如图，在 \( \triangle ABC \) 中，\( AB=BC=2 \)，\( \angle ABC=120^\circ \). 若平面 \(ABC\) 外的点 \(P\) 和线段 \(AC\) 上的点 \(D\)，满足 \( PD=DA \)，\( PB=BA \)，则四面体 \(PBCD\) 的体积的最大值是\fillinblank{}.

\bitmapfigure[max width=0.42\textwidth]{img/2016/zhejiang_science/q14_fig1.png}
\end{problem}

\begin{answer}
\(\dfrac12\)
\end{answer}

\begin{solution}
由余弦定理，\(AC^2=2^2+2^2-2\cdot2\cdot2\cos120^\circ=12\)，所以 \(AC=2\sqrt3\). 取平面 \(ABC\) 为 \(z=0\)，令
\(A=(0,0,0)\)，\(C=(2\sqrt3,0,0)\)，\(B=(\sqrt3,1,0)\).
这些坐标满足 \(AB=BC=2\) 且 \(\angle ABC=120^\circ\). 设 \(AD=x\)，其中 \(0\leqslant x\leqslant2\sqrt3\)，则 \(D=(x,0,0)\). 设 \(P\) 在平面 \(ABC\) 上的投影为 \(H=(u,v,0)\)，且 \(PH=h\). 由 \(PB=2\)、\(PD=x\)，有
\[
(u-\sqrt3)^2+(v-1)^2+h^2=4
\]
\[
(u-x)^2+v^2+h^2=x^2
\]
两式相减得
\[
v=(x-\sqrt3)u
\]
所以 \(H\) 在线 \(v=(x-\sqrt3)u\) 上. 因为
\[
h^2=x^2-DH^2
\]
要使 \(h\) 最大，应使 \(DH\) 为点 \(D=(x,0)\) 到该直线的距离. 因而
\(DH_{\min}=\frac{x|x-\sqrt3|}{\sqrt{1+(x-\sqrt3)^2}}\)，\(h_{\max}=\frac{x}{\sqrt{1+(x-\sqrt3)^2}}\).
又 \(B\) 到直线 \(AC\) 的距离为 \(1\)，故
\[
[BCD]=\frac12(2\sqrt3-x)
\]
从而
\[
V_{\max}(x)=\frac{1}{3}[BCD]h_{\max}
=\frac{x(2\sqrt3-x)}{6\sqrt{1+(x-\sqrt3)^2}}
\]
令 \(t=x-\sqrt3\)，则 \(t\in[-\sqrt3,\sqrt3]\)，且
\[
V_{\max}(x)=\frac{3-t^2}{6\sqrt{1+t^2}}
\]
令 \(u=t^2\)，则 \(0\leqslant u\leqslant3\)，并且
\[
\frac{\mathrm d}{\mathrm du}\left(\frac{3-u}{6\sqrt{1+u}}\right)
=-\frac{5+u}{12(1+u)^{3/2}}<0
\]
所以当 \(u=0\)，即 \(x=\sqrt3\) 时体积最大，最大值为
\[
V_{\max}=\frac{3}{6}=\frac12
\]
\end{solution}

\begin{problem}
已知平面向量 \(\bs a\)，\(\bs b\)，\( |\bs a|=1 \)，\( |\bs b|=2 \)，若对任意单位向量 \( \bs e \)，均有 \( |\bs a\cdot \bs e|+|\bs b\cdot \bs e|\le \sqrt6 \)，则 \( \bs a\cdot \bs b \) 的最大值是\fillinblank{}.
\end{problem}

\begin{answer}
\(\dfrac12\)
\end{answer}

\begin{solution}
对任意单位向量 \(\bs e\)，有
\[
|\bs a\cdot\bs e|+|\bs b\cdot\bs e|
=\max_{\varepsilon,\delta\in\{1,-1\}}
(\varepsilon\bs a+\delta\bs b)\cdot\bs e
\]
再对单位向量 \(\bs e\) 取最大值，得到
\[
\max_{|\bs e|=1}\bigl(|\bs a\cdot\bs e|+|\bs b\cdot\bs e|\bigr)
=\max\{|\bs a+\bs b|,|\bs a-\bs b|\}
=\sqrt{5+2|\bs a\cdot\bs b|}
\]
题设不等式因此等价于
\[
\sqrt{5+2|\bs a\cdot\bs b|}\leqslant\sqrt6
\]
即 \(|\bs a\cdot\bs b|\leqslant\dfrac12\). 所以 \(\bs a\cdot\bs b\leqslant\dfrac12\). 取夹角满足 \(\cos\theta=\dfrac14\) 的向量 \(\bs a\)，\(\bs b\)，则 \(\bs a\cdot\bs b=2\cos\theta=\dfrac12\)，且题设条件取等号，故最大值为 \(\dfrac12\).
\end{solution}

\section{解答题}

\begin{problem}
在 \( \triangle ABC \) 中，内角 \(A\)，\(B\)，\(C\) 所对的边分别为 \(a\)，\(b\)，\(c\)，已知 \( b+c=2a\cos B \).

\begin{enumerate}
\item 证明：\( A=2B \)；
\item 若 \( \triangle ABC \) 的面积 \( S=\dfrac{a^2}{4} \)，求角 \(A\) 的大小.
\end{enumerate}
\end{problem}

\begin{solution}
\begin{enumerate}
\item 由正弦定理，
\(\frac{b}{a}=\frac{\sin B}{\sin A}\)，\(\frac{c}{a}=\frac{\sin C}{\sin A}\).
将已知条件 \(b+c=2a\cos B\) 两边同除以 \(a\)，得
\[
\sin B+\sin C=2\sin A\cos B
\]
又 \(C=\pi-A-B\)，所以 \(\sin C=\sin(A+B)\). 因而
\[
\sin B+\sin(A+B)=2\sin A\cos B
\]
即
\[
\sin B=\sin A\cos B-\cos A\sin B=\sin(A-B)
\]
由于 \(\sin B>0\)，等式 \(\sin B=\sin(A-B)\) 迫使 \(A-B\in(0,\pi)\). 在 \(0<B<\pi\)、\(0<A-B<\pi\) 下，由正弦值相等可得
\[
A-B=B\quad\text{或}\quad A-B=\pi-B
\]
第二种情形给出 \(A=\pi\)，不符合三角形内角条件，所以 \(A=2B\).

\item 由面积公式和正弦定理，
\[
\frac12bc\sin A
=\frac12\cdot\frac{a\sin B}{\sin A}\cdot\frac{a\sin C}{\sin A}\cdot\sin A
=\frac{a^2\sin B\sin C}{2\sin A}
\]
结合 \(S=\dfrac{a^2}{4}\)，得
\[
2\sin B\sin C=\sin A
\]
由第一问 \(A=2B\)，且 \(\sin A=2\sin B\cos B\)，所以
\[
\sin C=\cos B
\]
又 \(C=\pi-3B\)，且 \(B\)，\(C\in(0,\pi)\). 因为 \(0<B<\pi/3\)，所以 \(\cos B>0\)，从
\(\sin C=\sin\left(\dfrac\pi2-B\right)\) 得
\[
C=\frac\pi2-B\quad\text{或}\quad C=\frac\pi2+B
\]
分别代入 \(C=\pi-3B\)，得到 \(B=\dfrac\pi4\) 或 \(B=\dfrac\pi8\). 因此
\[
A=2B=\frac\pi2\quad\text{或}\quad A=\frac\pi4
\]
\end{enumerate}
\end{solution}

\begin{problem}
如图，在三棱台 \( ABC-DEF \) 中，已知平面 \( BCFE\perp \) 平面 \(ABC\)，\( \angle ACB=90^\circ \)，\( BE=EF=FC=1 \)，\( BC=2 \)，\( AC=3 \).

\begin{enumerate}
\item 求证：\( BF\perp \) 平面 \(ACFD\)；
\item 求二面角 \( B-AD-F \) 的平面角的余弦值.
\end{enumerate}

\bitmapfigure[max width=0.42\textwidth]{img/2016/zhejiang_science/q17_fig1.png}
\end{problem}

\begin{solution}
\begin{enumerate}
\item 延长 \(AD\)、\(BE\)、\(CF\)，设它们交于点 \(K\). 由于 \(EF\parallel BC\)，在三角形 \(KBC\) 中有
\[
\frac{EF}{BC}=\frac12=\frac{KE}{KB}=\frac{KF}{KC}
\]
又 \(BE=FC=1\)，所以 \(KB=KC=2\)，从而 \(\triangle BCK\) 为边长 \(2\) 的等边三角形，且 \(F\) 是 \(CK\) 的中点.

因为平面 \(BCFE\perp\) 平面 \(ABC\)，且 \(AC\subset\) 平面 \(ABC\)、\(AC\perp BC\)，所以 \(AC\perp\) 平面 \(BCFE\)，从而 \(AC\perp BF\). 在等边三角形 \(BCK\) 中，中线 \(BF\) 同时是高，故 \(BF\perp CK\)，即 \(BF\perp CF\). 直线 \(AC\) 与 \(CF\) 在点 \(C\) 相交，且均在平面 \(ACFD\) 内，所以
\[
BF\perp\text{平面 }ACFD
\]

\item 由 \(DF\parallel AC\) 且 \(EF/BC=1/2\)，点 \(D\) 是 \(AK\) 的中点. 因而平面 \(BAD\) 就是平面 \(ABK\)，平面 \(FAD\) 就是平面 \(ACK\).

取 \(O\) 为 \(BC\) 的中点，以 \(OB\) 方向为 \(x\) 轴，以 \(OK\) 方向为 \(z\) 轴，建立空间直角坐标系，使平面 \(ABC\) 为 \(z=0\). 由 \(\triangle BCK\) 为边长 \(2\) 的等边三角形以及 \(AC\perp BC\)、\(AC=3\)，可取
\(B=(1,0,0)\)，\(C=(-1,0,0)\)，\(K=(0,0,\sqrt3)\)，\(A=(-1,-3,0)\).
于是
\(\overrightarrow{AB}=(2,3,0)\)，\(\overrightarrow{AK}=(1,3,\sqrt3)\)，\(\overrightarrow{AC}=(0,3,0)\).
平面 \(ABK\) 和平面 \(ACK\) 的一个法向量分别为
\[
\bs n_1=\overrightarrow{AB}\times\overrightarrow{AK}
=\sqrt3(3,-2,\sqrt3)
\]
\[
\bs n_2=\overrightarrow{AC}\times\overrightarrow{AK}
=3(\sqrt3,0,-1)
\]
所以二面角 \(B-AD-F\) 的平面角 \(\theta\) 满足
\[
\cos\theta
=\frac{|(3,-2,\sqrt3)\cdot(\sqrt3,0,-1)|}{4\cdot2}
=\frac{2\sqrt3}{4\cdot2}
=\frac{\sqrt3}{4}
\]
\end{enumerate}
\end{solution}

\begin{problem}
已知 \( a\ge 3 \)，函数 \( F(x)=\min\{2|x-1|,x^2-2ax+4a-2\} \)，其中 \( \min(p,q)=\begin{cases}
p, & p\le q,\\
q, & p>q.
\end{cases} \)

\begin{enumerate}
\item 求使得等式 \( F(x)=x^2-2ax+4a-2 \) 成立的 \( x \) 的取值范围；
\item\begin{enumerate}
\item 求 \( F(x) \) 的最小值 \( m(a) \)；
\item 求 \( F(x) \) 在 \( [0,6] \) 上的最大值 \( M(a) \).
\end{enumerate}
\end{enumerate}
\end{problem}

\begin{solution}
\begin{enumerate}
\item 令 \(f(x)=2|x-1|\)，\(g(x)=x^2-2ax+4a-2\). 由 \(F(x)=g(x)\) 等价于 \(g(x)\leqslant f(x)\)，分情况讨论.

当 \(x\leqslant1\) 时，
\[
g(x)-f(x)=x^2-2ax+4a-2-2(1-x)
=x^2+2(a-1)(2-x)>0
\]
其中 \(a\geqslant3\) 且 \(2-x\geqslant1\). 因而此时 \(g(x)>f(x)\)，不满足条件.

当 \(x>1\) 时，
\[
g(x)-f(x)=x^2-2ax+4a-2-2(x-1)
=(x-2)(x-2a)
\]
所以 \(g(x)\leqslant f(x)\) 等价于 \(2\leqslant x\leqslant2a\). 综上，所求范围为
\[
x\in[2,2a]
\]

\item
\begin{enumerate}
\item 对任意 \(x\)，有
\(f(x)\geqslant0=f(1)\)，\(g(x)=(x-a)^2-a^2+4a-2\geqslant-a^2+4a-2=g(a)\).
由于 \(F(x)=\min\{f(x),g(x)\}\)，故
\[
m(a)=\min\{0,-a^2+4a-2\}
\]
当 \(3\leqslant a\leqslant2+\sqrt2\) 时，\(-a^2+4a-2\geqslant0\)；当 \(a>2+\sqrt2\) 时，\(-a^2+4a-2<0\). 因此
\[
m(a)=
\begin{cases}
0, & 3\leqslant a\leqslant2+\sqrt2,\\
-a^2+4a-2, & a>2+\sqrt2.
\end{cases}
\]

\item 当 \(0\leqslant x\leqslant2\) 时，由第一问可知 \(g(x)\geqslant f(x)\)，所以 \(F(x)=f(x)\)，从而
\[
F(x)\leqslant\max\{f(0),f(2)\}=2
\]
当 \(2\leqslant x\leqslant6\) 时，由 \(a\geqslant3\) 得 \(x\leqslant6\leqslant2a\)，所以 \(F(x)=g(x)\). 二次函数 \(g(x)\) 开口向上，故在 \([2,6]\) 上的最大值在端点取得，且
\(g(2)=2\)，\(g(6)=34-8a\).
所以
\[
M(a)=
\begin{cases}
34-8a, & 3\leqslant a<4,\\
2, & a\geqslant4.
\end{cases}
\]
\end{enumerate}
\end{enumerate}
\end{solution}

\begin{problem}
如图，设椭圆 \( C:\dfrac{x^2}{a^2}+y^2=1 \)（\( a>1 \)）.

\begin{enumerate}
\item 求直线 \( y=kx+1 \) 被椭圆截得的线段长（用 \(a\)，\(k\) 表示）；
\item 若任意以点 \( A(0,1) \) 为圆心的圆与椭圆至多有三个公共点，求椭圆的离心率的取值范围.
\end{enumerate}

\bitmapfigure[max width=0.42\textwidth]{img/2016/zhejiang_science/q19_fig1.png}
\end{problem}

\begin{solution}
\begin{enumerate}
\item 将直线 \(y=kx+1\) 代入椭圆方程，得
\[
\frac{x^2}{a^2}+(kx+1)^2=1
\]
即
\[
(1+a^2k^2)x^2+2a^2kx=0
\]
设两个交点的横坐标为 \(x_1\)，\(x_2\)，则
\(x_1=0\)，\(x_2=-\frac{2a^2k}{1+a^2k^2}\).
直线斜率为 \(k\)，所以弦长为
\[
\sqrt{1+k^2}\,|x_1-x_2|
=\frac{2a^2|k|\sqrt{1+k^2}}{1+a^2k^2}
\]

\item 记 \(A=(0,1)\). 对任意 \(k>0\)，直线 \(y=kx+1\) 与椭圆除 \(A\) 外的另一个交点在 \(y\) 轴左侧；其到 \(A\) 的距离就是第一问所得弦长，记为 \(L(k)\). 椭圆关于 \(y\) 轴对称，因此若存在 \(k_1\ne k_2>0\) 使 \(L(k_1)=L(k_2)\)，则左侧有两个不同的点到 \(A\) 等距，关于 \(y\) 轴的两个对称点也到 \(A\) 等距，于是以 \(A\) 为圆心的相应圆有四个公共点. 所以题设条件等价于 \(L(k)\) 在 \(k>0\) 上为单调函数.

令 \(t=k^2>0\)，由第一问有
\[
\frac{L(k)^2}{4}
=\Phi(t)
=\frac{a^4t(1+t)}{(1+a^2t)^2}
\]
求导得
\[
\Phi'(t)=\frac{a^4\left[1+(2-a^2)t\right]}{(1+a^2t)^3}
\]
当 \(a^2\leqslant2\) 时，\(\Phi'(t)>0\)，所以 \(L(k)\) 在 \(k>0\) 上严格递增，任意以 \(A\) 为圆心的圆至多与椭圆有三个公共点.

更具体地，当 \(a^2\leqslant2\) 时，对任意有限的 \(t>0\)，有
\[
1-\Phi(t)=\frac{1+a^2(2-a^2)t}{(1+a^2t)^2}>0
\]
故 \(0<L(k)<2\). 除点 \(A\) 外，左、右两侧等距点至多各有一个；圆心为 \(A\) 的半径为 \(2\) 的圆还只包含椭圆的下顶点，因而公共点数不超过 \(3\).

当 \(a^2>2\) 时，\(\Phi'(t)\) 在 \(t=\dfrac1{a^2-2}\) 前为正、后为负，故 \(\Phi\) 先增后减，存在两个不同的正数 \(t_1\)，\(t_2\) 使 \(\Phi(t_1)=\Phi(t_2)\)，从而会产生四个公共点. 因此题设条件等价于
\[
1<a\leqslant\sqrt2
\]
椭圆的离心率为
\[
e=\frac{\sqrt{a^2-1}}{a}=\sqrt{1-\frac1{a^2}}
\]
所以
\[
0<e\leqslant\frac{\sqrt2}{2}
\]
\end{enumerate}
\end{solution}

\begin{problem}
设数列满足 \( \left|a_n-\dfrac{a_{n+1}}{2}\right|\le 1 \)，\( n\in\N^* \).

\begin{enumerate}
\item 求证：\( |a_n|\ge 2^{n-1}\left( |a_1|-2 \right) \)，\( n\in\N^* \)；
\item 若 \( |a_n|\le \left( \dfrac32 \right)^n \)，\( n\in\N^* \)，证明：\( |a_n|\le 2 \)，\( n\in\N^* \).
\end{enumerate}
\end{problem}

\begin{solution}
\begin{enumerate}
\item 由题设不等式及反三角不等式，
\[
\frac{|a_{n+1}|}{2}\geqslant |a_n|-\left|a_n-\frac{a_{n+1}}2\right|\geqslant |a_n|-1
\]
从而
\[
|a_{n+1}|-2\geqslant2(|a_n|-2)
\]
反复使用上式，得到
\[
|a_n|-2\geqslant2^{n-1}(|a_1|-2)
\]
于是
\(|a_n|\geqslant2+2^{n-1}(|a_1|-2) \geqslant2^{n-1}(|a_1|-2)\)，\(n\in\N^*\).

\item 固定任意 \(n\in\N^*\). 若 \(|a_n|>2\)，则对任意 \(m\geqslant n\)，由同一递推关系迭代可得
\[
|a_m|-2\geqslant2^{m-n}(|a_n|-2)
\]
另一方面，由题设上界，
\[
|a_m|-2<|a_m|\leqslant\left(\frac32\right)^m
\]
合并两式，得到
\[
0<|a_n|-2\leqslant2^{n-m}\left(\frac32\right)^m
=2^n\left(\frac34\right)^m
\]
令 \(m\to\infty\)，右端趋于 \(0\)，与 \(|a_n|-2>0\) 矛盾. 因此不存在 \(|a_n|>2\) 的 \(n\)，即
\(|a_n|\leqslant2\)，\(n\in\N^*\).
\end{enumerate}
\end{solution}
